%============================================================
% DISCUSSION
%============================================================

\section{Discussion}
\label{sec:discussion}

The aim of this study was to investigate heterogeneity in baseline
headache frequency and to assess whether a Poisson finite mixture
model could provide a more informative representation of the patient
population than a single Poisson distribution. The analysis provided
evidence of substantial variation in headache frequency and supported
the presence of more than one underlying frequency level among the
patients.

%============================================================
\subsection{Heterogeneity in Headache Frequency}
%============================================================

The initial descriptive analysis showed considerable variability in
the number of headache days reported by patients. The variance of
44.98 was substantially larger than the mean of 16.18, giving a
variance-to-mean ratio of 2.78. This is inconsistent with the
equidispersion property of a standard Poisson distribution.

The over-dispersion observed in these data may reflect differences
between patients in their underlying susceptibility to headaches.
Rather than treating the additional variation only as a statistical
problem, the finite mixture approach allowed the analysis to
investigate whether the population could be represented by several
latent groups with different headache-frequency rates.

This interpretation is consistent with the purpose of finite mixture
models, which allow an observed population to be represented as a
combination of unobserved subpopulations with different parameter
values \citep{McLachlanPeel2000}.

%============================================================
\subsection{Evidence from the NPMLE Analysis}
%============================================================

The NPMLE analysis provided an exploratory assessment of the
underlying mixing distribution. The gradient function showed
departures from a homogeneous one-component distribution and
suggested that more than one support point may be present.

The NPMLE was particularly useful because it allowed the underlying
mixing distribution to be explored before fixing the number of
components in the finite mixture models. This is consistent with the
role of NPMLE in mixture modelling, where the mixing distribution can
be estimated without imposing a specific finite number of support
points in advance \citep{Lindsay1983,Lindsay1983b}.

The gradient analysis should, however, be interpreted as exploratory
evidence rather than as a definitive test for the exact number of
components. The final choice of the finite mixture structure was
therefore based on the fitted models, information criteria, and the
interpretability of the estimated groups.

%============================================================
\subsection{Choice of the Number of Components}
%============================================================

The comparison of the candidate models demonstrated the importance
of considering both model fit and model complexity. The single
Poisson model performed substantially worse than the finite mixture
models, providing strong evidence that a homogeneous Poisson model
was inadequate for these data.

Among the mixture models, AIC reached its minimum for the
three-component model, whereas BIC reached its minimum for the
two-component model. This difference reflects the stronger penalty
for model complexity imposed by BIC \citep{Akaike1974,Schwarz1978}.

The results therefore do not provide unanimous statistical evidence
for exactly three components. Nevertheless, the three-component
model was retained as the principal model for interpretation because
it provided the lowest AIC and produced three clearly separated and
interpretable headache-frequency levels.

This decision should be viewed as a modelling choice rather than as
evidence that the population necessarily contains exactly three
biologically distinct groups. The latent components are statistical
representations of heterogeneity in the observed distribution.

%============================================================
\subsection{Interpretation of the Three Latent Groups}
%============================================================

The estimated component means from the three-component model were
9.61, 13.80, and 22.95 headache days. These estimates indicate
substantial differences in the underlying headache-frequency rates.

The lowest-frequency component represented approximately one quarter
of the sample, while the intermediate- and high-frequency components
each represented approximately two-fifths of the population. The
high-frequency group therefore represented a substantial proportion
of the patients rather than a small group of extreme observations.

The posterior classification plot also showed a clear separation of
patients across the three estimated frequency levels. This supports
the interpretation that the mixture model is capturing meaningful
heterogeneity in the observed distribution.

However, the latent groups should not be interpreted as fixed
clinical categories without additional clinical information. They
are statistical groups defined by differences in the estimated
headache-frequency distributions.

%============================================================
\subsection{Agreement Between CAMAN and FlexMix}
%============================================================

The analysis used two different approaches to investigate the
mixture structure. CAMAN was used for the NPMLE analysis, while
FlexMix was used to estimate and classify patients under the
three-component finite mixture model.

The estimated component means obtained from the two approaches were
similar, particularly for the highest-frequency component. Although
the estimated mixing proportions differed to some extent, both
methods identified a comparable set of low-, intermediate-, and
high-frequency rates.

This agreement provides additional support for the existence of
heterogeneity in the headache-frequency distribution. Differences
between the estimated proportions are not unexpected because the
methods use different estimation procedures and because mixture
components are latent rather than directly observed.

It is also important to recognize that component labels are arbitrary
in finite mixture models. Therefore, Component 1 in one software
implementation does not necessarily correspond to Component 1 in
another implementation. Comparisons should instead be based on the
estimated component parameters and their ordering.

%============================================================
\subsection{Model Adequacy and the Upper Boundary}
%============================================================

Although the three-component model captured the major structure of
the observed distribution, the model diagnostics revealed an
important limitation. The number of patients reporting the maximum
observed value of 28 headache days was substantially larger than
expected under the fitted mixture.

Thirty-eight patients were observed at 28 headache days, compared
with approximately 6.79 expected under the three-component model.
The corresponding Pearson residual was 11.98.

This concentration at the upper boundary suggests that the observed
distribution contains a feature that is not fully represented by the
ordinary Poisson mixture. One possible explanation is that the
four-week observation period imposes a natural upper limit of
28 headache days. Patients who experience headaches throughout the
entire observation period are therefore concentrated at the boundary
rather than continuing along an unrestricted count distribution.

This feature is important because the standard Poisson mixture treats
the count distribution as though its support were not constrained by
such a maximum. A model that explicitly accounts for a bounded or
zero/upper-inflated outcome could potentially provide a better
description of this aspect of the data.

The finding does not invalidate the mixture analysis. Instead, it
shows that the finite mixture model captures the broad latent
heterogeneity while leaving a specific boundary feature unexplained.

%============================================================
\subsection{Effect of Age}
%============================================================

The age-adjusted three-component model provided evidence that the
relationship between age and headache frequency was not uniform
across the latent population.

Age was statistically significant in one of the SAS mixture
components, with an estimated coefficient of $-0.02414$. On the
original rate scale, this corresponds to a multiplicative factor of
approximately $0.976$ for a one-year increase in age. Thus, within
that component, the expected headache frequency decreased by
approximately 2.4\% for each additional year of age.

The age effects in the other two components were not statistically
significant. This suggests that the association between age and
headache frequency may differ across latent groups rather than
following a single population-wide effect.

The component-specific nature of this result is important. It should
not be interpreted as evidence that increasing age generally reduces
headache frequency among all patients. Rather, the result indicates
that age was associated with the expected headache rate within one
particular latent component.

%============================================================
\subsection{Strengths of the Analysis}
%============================================================

A major strength of the analysis was the use of several complementary
methods rather than relying on a single fitted model. The descriptive
analysis identified over-dispersion, the NPMLE provided an
exploratory assessment of the mixing distribution, and finite
mixture models were then compared using information criteria.

The use of both CAMAN and FlexMix also provided an opportunity to
examine whether similar latent structures could be obtained using
different estimation approaches. The broadly similar component means
obtained from the two methods strengthen the interpretation of the
identified heterogeneity.

The analysis was also implemented in both R and SAS. This provided
an additional check on the model estimates and allowed the age-
adjusted mixture model to be examined using SAS.

%============================================================
\subsection{Limitations}
%============================================================

Several limitations should be considered when interpreting the
results.

First, the finite mixture components are latent statistical
subpopulations and should not automatically be interpreted as
distinct clinical disease categories. Additional clinical variables
would be required to establish whether the identified groups
correspond to meaningful clinical phenotypes.

Second, the AIC and BIC results did not select the same number of
components. AIC favored three components, whereas BIC favored two.
The three-component model was therefore selected as the principal
model for interpretation, but the two-component solution remains a
plausible alternative.

Third, the model showed substantial lack of fit at the upper
boundary of 28 headache days. This indicates that the ordinary
Poisson mixture does not fully account for the bounded nature of the
outcome.

Finally, the analysis focused primarily on baseline headache
frequency and age. Other patient characteristics may also contribute
to the observed heterogeneity and could be considered in future
models.

%============================================================
\subsection{Implications for Further Analysis}
%============================================================

The results suggest several directions for further statistical
investigation. A model that explicitly accounts for the upper
boundary of 28 headache days could be considered to address the
concentration of observations at the maximum value.

Additional covariates could also be incorporated into the mixture
model to determine whether demographic or clinical characteristics
help explain membership in the latent frequency groups.

Future work could further compare alternative count distributions and
mixture specifications using both statistical fit and substantive
interpretability. Such analyses would help determine whether the
three-component Poisson mixture provides the most useful description
of the underlying patient heterogeneity or whether a more flexible
model is warranted.

%============================================================
\subsection{Overall Interpretation}
%============================================================

Taken together, the findings indicate that baseline headache
frequency was heterogeneous across the study population and that a
single Poisson distribution was insufficient to describe the
observed variation.

The finite mixture analysis provided evidence for several underlying
headache-frequency levels, with the three-component model identifying
low-, intermediate-, and high-frequency groups. Although the
three-component model was not preferred by every information
criterion, it provided the lowest AIC and produced an interpretable
latent structure.

At the same time, the substantial lack of fit at 28 headache days
shows that the mixture model does not explain every feature of the
observed distribution. The findings therefore support the use of
finite mixture modelling as a useful representation of heterogeneity
in these data, while also indicating areas where a more specialized
count model may be required.